The leading margin survives at the named positive $q$.  The proof below is
deliberately quantitative: it keeps the finite temporal cutoff, reflection,
initial mismatch, varying derivative packet, and nonzero source mass.

\begin{lemma}[Uniform finite-$q$ correction sign; proved here]
\label{lem:terminal-modal-finite-correction-sign}
If $m\ge64$, then for every $2\le n\le m$ and every shared coordinate
$0\le j\le n-2$,
\begin{equation}
 K_n(j)>0.
 \label{eq:terminal-modal-finite-correction-sign}
\end{equation}
More precisely, for $n\ge6$,
\begin{equation}
 \frac{K_n(j)}{q^3(1-q)^{n-2}}
 \ge \frac{179}{14400}.
 \label{eq:terminal-modal-finite-correction-margin}
\end{equation}
Consequently the shared sign \eqref{eq:terminal-modal-shared-sign} holds and
the proper-prefix chronology in
\cref{lem:terminal-modal-admission-chronology} is unconditional in this
parameter range.
\end{lemma}

\begin{proof}
Put $\delta=1-q$ and $\wp_n=P_n/q^2$.  Rescale the correction by
\begin{equation}
 \bar c_1=\delta\frac{C_1}{q^3},\qquad
 \bar c_n=\delta^{-(n-2)}\frac{C_n}{q^3}\quad(n\ge2).
 \label{eq:terminal-modal-finite-rescaling}
\end{equation}
Thus
\begin{equation}
 \frac{K_n}{q^3}=\delta^{n-2}
 \left(\bar c_n-\frac12\bar c_{n-1}\right).
 \label{eq:terminal-modal-finite-rescaled-target}
\end{equation}
After even reflection, the homogeneous recurrence for $\bar c_n$ is exactly
the $q=0$ recurrence in
\eqref{eq:terminal-modal-leading-correction-recurrence}.  Only its moving
frontier inputs change.  Indeed, with
\begin{align}
 E_n&=\delta^{-(n-2)}
 \left(\frac15-\frac\eta2\wp_{n-1}\right),
 \label{eq:terminal-modal-finite-E}\\
 F_n&=\frac{\delta^{-(n-2)}}4
 \left(\frac15-\frac\eta2\wp_{n-1}+\wp_n\right)
 -\frac1{5\delta^{n-1}},
 \label{eq:terminal-modal-finite-F}
\end{align}
the exact recurrence is
\begin{align}
 \bar z_n(j)&=2\bar c_n(j)-\bar c_{n-1}(j),&&0\le j\le n-2,\notag\\
 \bar z_n(n-1)&=2\bar c_n(n-1)+E_n,\notag\\
 \bar c_{n+1}|_{[0,n-1]}&=B_0\bar z_n,\qquad
 \bar c_{n+1}(n)=\frac12\bar c_n(n-1)+F_n,
 \label{eq:terminal-modal-finite-rescaled-recurrence}
\end{align}
where $B_0$ is the reflected $q=0$ stencil.  This follows directly by
substituting \eqref{eq:terminal-modal-finite-rescaling} in
\eqref{eq:terminal-modal-correction-shared-z}--
\eqref{eq:terminal-modal-correction-new-row}; the factor
$(1-q^2)/(1+q)=\delta$ is exactly what removes all positive-$q$ damping from
the homogeneous part.

The $q=0$ inputs are $E_n^\circ=-3/10$ and $F_n^\circ=9/40$.  Write
\begin{equation}
 \varepsilon_n=E_n+\frac3{10},\qquad h_n=\frac{\varepsilon_n}{4},
 \qquad \mu_n=F_n+\frac34E_n.
 \label{eq:terminal-modal-finite-source-split}
\end{equation}
The oriented source for the difference from the $q=0$ trajectory is
therefore
\begin{equation}
 h_n(1,2,-3)+\mu_n(0,0,1)
 \quad\hbox{on positions }n-2,n-1,n,
 \label{eq:terminal-modal-finite-source-packets}
\end{equation}
together with its reflection, for every $n\ge2$.  At $n=2$ the first
position is the seed and the two reflected copies coincide, giving the
required doubled seed entry.  Thus the even-line formula also includes the
two early half-line stencils that overlap the reflected boundary.

We first bound the scalar coefficients in this split.  For $r=n-2$ and
$t=\chi^{n-1}$, put
\begin{equation}
 A=\delta^{-r},\qquad B=A\sqrt t,qquad
 e_0(t)=\frac{t^2-10t+3}{10(1+t^2)}.
\end{equation}
Direct simplification of \eqref{eq:terminal-modal-finite-E} using
\eqref{eq:terminal-modal-frontier-optimum} gives
\begin{equation}
 \varepsilon_n
 =B\left(t^{-1/2}e_0(t)+\frac3{10}\right)
 +\frac3{10}(1-B)-\frac{Aq}{10}.
 \label{eq:terminal-modal-finite-epsilon-identity}
\end{equation}
Bernoulli's inequality gives $A\le16/15$, while
\begin{equation}
 1-q\le B=\sqrt\chi(1-q^2)^{-r/2}\le1.
 \label{eq:terminal-modal-finite-AB-bounds}
\end{equation}
The upper bound follows from $(1-q^2)^r\ge\chi$, using $rq\le1/16$;
the lower bound follows from $\sqrt\chi\ge1-q$.

Since $t\ge7/8$, $y=\sqrt t\ge14/15$.  Exact polynomial clearing gives
\begin{equation}
 -\frac1{375}\le t^{-1/2}e_0(t)+\frac3{10}\le0.
 \label{eq:terminal-modal-finite-e0-window}
\end{equation}
For the upper inequality the numerator factors as
$(y-1)(3y^4+4y^3+4y^2-6y-3)$, whose second factor is positive on
$[14/15,1]$.  For the lower inequality, after clearing the positive
denominator, the polynomial is
\[5(227y^5+75y^4-750y^2+227y+225).
\]
Its degree-five Bernstein coefficients on
$y=(14+z)/15$, $0\le z\le1$, are
\[
 \frac{924673}{151875},\quad \frac{15872}{10125},\quad
 \frac{238}{675},\quad \frac{25}{9},\quad \frac{46}{5},\quad20,
\]
all positive.  Because $q\le1/1024$, equations
\eqref{eq:terminal-modal-finite-epsilon-identity}--
\eqref{eq:terminal-modal-finite-e0-window} yield
\begin{equation}
 -\frac1{360}\le\varepsilon_n\le\frac{3q}{10},\qquad
 |\varepsilon_n|\le\frac1{360}.
 \label{eq:terminal-modal-finite-epsilon-pointwise}
\end{equation}
For example, the lower loss is at most
$1/375+(16/15)(1/1024)/10=133/48000<1/360$.

The same identity gives a total-variation bound.  Extend $r$ continuously,
keep $t=\chi^{r+1}$, and put
$\lambda=-\log(1-q)/q$ and $\kappa=-\log\chi/q$.  Then
\begin{equation}
 \left|\frac{dE}{dr}\right|
 =Aq\left|\lambda(e_0-q/10)-\kappa t e_0'\right|.
\end{equation}
On $7/8\le t\le1$, direct factorization gives
\begin{equation}
 |e_0|\le\frac3{10},\qquad |te_0'|\le\frac15,qquad
 |e_0-2te_0'|\le\frac1{10}.
\end{equation}
The first and third bounds reduce respectively to products of linear/cubic
factors of fixed sign.  For the middle bound, one cleared side is
$(t^2-3t-1)(t^2-2t-1)\ge0$, while the other is
$t^4+5t^3-5t+1>0$, increasing from its positive value at $7/8$.
Finally
\[
 0\le\lambda-1\le\frac q{1-q},qquad
 0\le\kappa-2\le\frac{2q^2}{1-q^2}.
\]
Consequently
\[
 \left|\lambda(e_0-q/10)-\kappa te_0'\right|
 \le\frac1{10}+\frac{2q}{5(1-q)}
 +\frac{2q^2}{5(1-q^2)}\le\frac18.
\]
Since the source interval has length at most $m$,
\begin{equation}
 \operatorname{TV}(\varepsilon_n)=\operatorname{TV}(E_n)
 \le \frac{16}{15}\frac{mq}{8}=\frac1{120},qquad
 |\varepsilon_{N-1}|+\operatorname{TV}(\varepsilon_n)\le\frac1{90}.
 \label{eq:terminal-modal-finite-epsilon-TV}
\end{equation}

For the mass packet, no later chronology result is needed.  Directly from
\eqref{eq:terminal-modal-finite-E}--\eqref{eq:terminal-modal-finite-F} and
the fixed-prefix frontier identity
\[
 \wp_n=s_n\left(\wp_{n-1}-\frac{2q}{5\eta}\right)
\]
one obtains, for every $n\ge2$,
\begin{equation}
 \mu_n=-\delta^{-(n-1)}q\,\mathfrak m_n,
 \qquad
 \mathfrak m_n=
 \frac{\delta\wp_{n-1}}4\frac{2\eta-s_n}{q}
 +\frac{\delta s_n+2\eta}{10\eta}.
 \label{eq:terminal-modal-finite-mu-M}
\end{equation}
This is a direct scalar calculation on the fixed prefix, so it does not
import the source-mass lemma proved later along the chronology.  Once that
chronology is licensed, comparison with
\eqref{eq:terminal-modal-source-mass} says merely that
$\mathfrak m_n=M_n/q^4$.  Formula
\eqref{eq:terminal-modal-finite-mu-M} already covers the early reflected
overlaps $n=2,3$; the focused exact verifier checks those two cases
separately.
If $s_n=f_{n-1}/f_n$ and $u=\chi^n$, then
\begin{equation}
 \frac{2\eta-s_n}{q}
 =\frac{2(1-u^2)/(1+u^2)-3q+q^3}{1-q^2}.
 \label{eq:terminal-modal-finite-mass-ratio}
\end{equation}
Using $u\ge7/8$, $\wp_{n-1}\le33/16$, and $s_n\le1$ in the exact mass
formula gives
\begin{align}
 0<\mathfrak m_n
 &=\frac{(1-q)\wp_{n-1}}4\frac{2\eta-s_n}{q}
 +\frac{(1-q)s_n+2\eta}{10\eta}\notag\\
 &\le\frac{990}{7232}+\frac25
 =\frac{9707}{18080}<\frac{43}{80}.
 \label{eq:terminal-modal-finite-mass-bound}
\end{align}
For positivity, the second term is at least $1/5$, whereas the first is at
least $-99q/(64(1+q))>-1/5$.  The upper bound uses
$2(1-u^2)/(1+u^2)\le30/113$.  Since
$\delta^{n-1}\ge1-(n-1)q\ge15/16$, we conclude
\begin{equation}
 -\frac{43}{75}q<\mu_n<0.
 \label{eq:terminal-modal-finite-mu-bound}
\end{equation}

It remains to control the stopped spatial response.  Let $L_k^\circ$ be the
$q=0$ kernel and write
$\ell_k(r)=[z^r]L_k^\circ$.  From
\eqref{eq:terminal-modal-correction-green},
\begin{equation}
 \ell_{1}(r)=\mathbf1_{\{r=0\}},\qquad
 \ell_k(r)=2^{-(k-1)}
 \sum_{a=\max\{|r|-1,0\}}^{k-2}\binom{k-2}{a}
 \quad(k\ge2).
 \label{eq:terminal-modal-finite-l-coefficients}
\end{equation}
Thus $0\le\ell_k(r)\le1/2$ for $k\ge2$.  For an oriented moving derivative
packet put
\begin{equation}
 w_k(d)=\ell_k(k-d)+2\ell_k(k-d-1)-3\ell_k(k-d-2),\qquad d\ge0.
\end{equation}
The binomial-tail formula shows that the part with $k\le d$ is nonnegative,
$w_{d+1}(d)=0$ apart from the harmless $(d,k)=(0,1)$ endpoint, and, for
$k\ge d+2$,
\begin{equation}
 w_k(d)=-\frac12\left\{\Pr[Z_{k-2}=d]
 +3\Pr[Z_{k-2}=d+1]\right\},\qquad
 Z_s\sim\operatorname{Bin}(s,1/2).
 \label{eq:terminal-modal-finite-derivative-negative-lobe}
\end{equation}
The two negative binomial sums are each two, because
$\sum_{s\ge h}\binom{s}{h}2^{-s}=2$.  Hence the negative lobe has total
$-4$.  A direct generating-function extraction, equivalently summing the
positive binomial atoms, gives
\begin{equation}
 \sum_{k\ge1}w_k(d)=-\frac83+\frac23(-1/2)^d,qquad
 \sum_{k\le d}w_k(d)+\mathbf1_{\{d=0\}}w_1(0)
 =\frac43+\frac23(-1/2)^d\le2.
 \label{eq:terminal-modal-finite-derivative-lobes}
\end{equation}

For a final shared coordinate $(N,j)$, the positive and reflected packets
have $d=N-2-j$ and $d'=N-2+j$, respectively, while $k=N-n$ and
$1\le k\le N-2$.  The first branch has a nonnegative lobe of mass at most
two followed by a negative lobe of mass four.  The reflected branch always
has $k\le d'$ and is a subinterval of its nonnegative lobe, also of mass at most
two.  Source chronology increases $n$ and hence traverses $k=N-n$ in
decreasing order, so it encounters any negative lobe before the nonnegative
lobe.  Therefore every stopped chronological prefix of the combined
reflected derivative response has value in $[-4,4]$.  Abel summation and
\eqref{eq:terminal-modal-finite-epsilon-TV} now give the uniform loss
\begin{equation}
 \left|\sum_{n=2}^{N-1}h_n
 \bigl(\hbox{reflected derivative response at $(N,j)$}\bigr)\right|
 \le4\left(|h_{N-1}|+\operatorname{TV}(h_n)\right)
 \le\frac1{90}.
 \label{eq:terminal-modal-finite-derivative-loss}
\end{equation}

The folded response to one newest-row point mass is at most one: for $k=1$
this is immediate, and for $k\ge2$ it is
$\ell_k(j-n)+\ell_k(j+n)\le1$.  Thus
\eqref{eq:terminal-modal-finite-mu-bound} bounds the total mass loss by
$43q(N-2)/75$.

The exact rescaled initial values are
\begin{equation}
 \bar c_1=-\frac\delta5,\qquad
 \bar c_2=\left(\frac25-\frac q2+\frac{q^2}{10},
 -\frac q4+\frac{q^2}{20}\right).
 \label{eq:terminal-modal-finite-initial-values}
\end{equation}
Write $h_1^{\rm base},h_2^{\rm base}$ for the perturbations relative to the
$q=0$ initial pair.  Their even-line $\ell_1$ norms are respectively $q/5$
and $q-q^2/5<q$.  The base response at time $N\ge3$ is
$L_{N-1}^\circ h_2^{\rm base}-B_0L_{N-2}^\circ h_1^{\rm base}$.  Every Laurent coefficient of
$L_k^\circ$ for $k\ge2$, and of $B_0L_1^\circ$, is at most $1/2$; convolution
by the Markov kernel $B_0$ cannot increase that maximum.  Hence the
homogeneous initial loss is at most $q/2+q/10=3q/5$.

For $N\ge6$, the leading proof above sharpens to
\begin{equation}
 d_N^\circ(j)\ge\frac{19}{320}.
 \label{eq:terminal-modal-leading-correction-late-margin}
\end{equation}
Indeed the seed and first-neighbor formulas are larger; the fixed-$j$
coefficients with time index at least two are at least $1/16$; and at time
index one
$h_{j,1}=(9j-7)/2^j$, which decreases for $j\ge4$ and gives $19/320$ at
$j=4$.  Since $q=1/(16m)\le1/(16N)$,
\[
 \frac{3q}{5}+\frac{43q(N-2)}{75}
 \le\frac{43N-41}{1200N}<\frac{43}{1200}.
\]
Combining this with \eqref{eq:terminal-modal-finite-derivative-loss} yields
\[
 \bar c_N(j)-\frac12\bar c_{N-1}(j)
 \ge\frac{19}{320}-\frac{43}{1200}-\frac1{90}
 =\frac{179}{14400},
\]
which proves \eqref{eq:terminal-modal-finite-correction-margin}.

Only $2\le N\le5$ remain.  At $N=2$ direct substitution in
\eqref{eq:terminal-modal-finite-initial-values} gives
$(\bar c_2-\bar c_1/2)(0)=(1-q)(5-q)/10>0$.  For $3\le N\le5$, the $q=0$
shared-coordinate margin is at least $1/80$.  The initial loss is at most
$3q/5$, and at most three mass packets cost at most $43q/25$.
For the derivative packet, direct substitution in
\eqref{eq:terminal-modal-finite-E} gives
\begin{equation}
 \varepsilon_2=\frac{q(1+10q-q^2)}{10(1+q^2)}<\frac q8.
\end{equation}
The continuous bound used in \eqref{eq:terminal-modal-finite-epsilon-TV}
gives $|d\varepsilon/dr|\le2q/15$.  Here the source interval has length at most
two, so its variation is at most $4q/15$ and
\[
 |\varepsilon_{N-1}|+\operatorname{TV}(\varepsilon_n)
 \le |\varepsilon_2|+2\operatorname{TV}(\varepsilon_n)
 <\frac{79}{120}q.
\]
The same Abel bound therefore makes the complete perturbation smaller than
\[
 \left(\frac35+\frac{43}{25}+\frac{79}{120}\right)q
 =\frac{1787}{600}q<3q\le\frac3{1024}<\frac1{80}.
\]
Thus the leading margin also remains positive at $3\le N\le5$.
This completes \eqref{eq:terminal-modal-finite-correction-sign}.  Finally
\eqref{eq:terminal-modal-prefix-sign-split} gives $D_n\ge0$, and
\cref{lem:terminal-modal-chronology-reduction} proves
\cref{lem:terminal-modal-admission-chronology}.
\end{proof}
